\documentclass[10pt,a4paper]{amsart}
\usepackage{amsmath,amssymb,amsthm,graphicx}
\usepackage{verbatim}
\newenvironment{reference}{\comment}{\endcomment}
%\newenvironment{reference}{}{}
\newenvironment{slogan}{\comment}{\endcomment}
\newenvironment{history}{\comment}{\endcomment}

% For commutative diagrams we use Xy-pic
\usepackage[all]{xy}

% We use 2cell for 2-commutative diagrams.
\xyoption{2cell}
% Theorem environments.
%
\theoremstyle{plain}
\newtheorem{theorem}[subsection]{Theorem}
\newtheorem{proposition}[subsection]{Proposition}
\newtheorem{lemma}[subsection]{Lemma}

\theoremstyle{definition}
\newtheorem{definition}[subsection]{Definition}
\newtheorem{example}[subsection]{Example}
\newtheorem{exercise}[subsection]{Exercise}
\newtheorem{situation}[subsection]{Situation}

\theoremstyle{remark}
\newtheorem{remark}[subsection]{Remark}
\newtheorem{remarks}[subsection]{Remarks}

\numberwithin{equation}{subsection}

% Macros
%
\def\lim{\mathop{\mathrm{lim}}\nolimits}
\def\colim{\mathop{\mathrm{colim}}\nolimits}
\def\Spec{\mathop{\mathrm{Spec}}}
\def\Hom{\mathop{\mathrm{Hom}}\nolimits}
\def\Ext{\mathop{\mathrm{Ext}}\nolimits}
\def\SheafHom{\mathop{\mathcal{H}\!\mathit{om}}\nolimits}
\def\SheafExt{\mathop{\mathcal{E}\!\mathit{xt}}\nolimits}
\def\Sch{\mathit{Sch}}
\def\Mor{\mathop{\mathrm{Mor}}\nolimits}
\def\Ob{\mathop{\mathrm{Ob}}\nolimits}
\def\Sh{\mathop{\mathit{Sh}}\nolimits}
\def\NL{\mathop{N\!L}\nolimits}
\def\CH{\mathop{\mathrm{CH}}\nolimits}
\def\proetale{{pro\text{-}\acute{e}tale}}
\def\etale{{\acute{e}tale}}
\def\QCoh{\mathit{QCoh}}
\def\Ker{\mathop{\mathrm{Ker}}}
\def\Im{\mathop{\mathrm{Im}}}
\def\Coker{\mathop{\mathrm{Coker}}}
\def\Coim{\mathop{\mathrm{Coim}}}

% Boxtimes
%
\DeclareMathSymbol{\boxtimes}{\mathbin}{AMSa}{"02}

%
% Macros for moduli stacks/spaces
%
\def\QCohstack{\mathcal{QC}\!\mathit{oh}}
\def\Cohstack{\mathcal{C}\!\mathit{oh}}
\def\Spacesstack{\mathcal{S}\!\mathit{paces}}
\def\Quotfunctor{\mathrm{Quot}}
\def\Hilbfunctor{\mathrm{Hilb}}
\def\Curvesstack{\mathcal{C}\!\mathit{urves}}
\def\Polarizedstack{\mathcal{P}\!\mathit{olarized}}
\def\Complexesstack{\mathcal{C}\!\mathit{omplexes}}
% \Pic is the operator that assigns to X its picard group, usage \Pic(X)
% \Picardstack_{X/B} denotes the Picard stack of X over B
% \Picardfunctor_{X/B} denotes the Picard functor of X over B
\def\Pic{\mathop{\mathrm{Pic}}\nolimits}
\def\Picardstack{\mathcal{P}\!\mathit{ic}}
\def\Picardfunctor{\mathrm{Pic}}
\def\Deformationcategory{\mathcal{D}\!\mathit{ef}}
\usepackage{cleveref}
\crefname{originallemma}{original lemma}{original lemmas}
\Crefname{originallemma}{Original lemma}{Original lemmas}
\makeatletter
\newlabel{algebra-lemma-Noetherian-complete-local-Nagata}{{032W}{source}{}{sourcepointer-1}{}}
\newlabel{algebra-lemma-Noetherian-complete-local-Nagata@cref}{{[originallemma][0][]032W}{[1][0][]source}{}{}{}}
\makeatother
\crefname{originallemma}{original lemma}{original lemmas}
\Crefname{originallemma}{Original lemma}{Original lemmas}
\makeatletter
\newlabel{algebra-lemma-descent-normal}{{033G}{source}{}{sourcepointer-2}{}}
\newlabel{algebra-lemma-descent-normal@cref}{{[originallemma][0][]033G}{[1][0][]source}{}{}{}}
\makeatother
\crefname{originallemma}{original lemma}{original lemmas}
\Crefname{originallemma}{Original lemma}{Original lemmas}
\makeatletter
\newlabel{more-algebra-lemma-application-formal-glueing}{{0ALK}{source}{}{sourcepointer-3}{}}
\newlabel{more-algebra-lemma-application-formal-glueing@cref}{{[originallemma][0][]0ALK}{[1][0][]source}{}{}{}}
\makeatother
\crefname{originalremark}{original remark}{original remarks}
\Crefname{originalremark}{Original remark}{Original remarks}
\makeatletter
\newlabel{more-algebra-remark-formal-glueing-algebras}{{05EU}{source}{}{sourcepointer-4}{}}
\newlabel{more-algebra-remark-formal-glueing-algebras@cref}{{[originalremark][0][]05EU}{[1][0][]source}{}{}{}}
\makeatother
\crefname{originallemma}{original lemma}{original lemmas}
\Crefname{originallemma}{Original lemma}{Original lemmas}
\makeatletter
\newlabel{morphisms-lemma-nagata-normalization}{{035S}{source}{}{sourcepointer-5}{}}
\newlabel{morphisms-lemma-nagata-normalization@cref}{{[originallemma][0][]035S}{[1][0][]source}{}{}{}}
\makeatother
\title{Normality away from the special fibre after completion forces the normalization to be finite and compatible with completion}
\author{The Stacks Project Authors}
\date{stacks-project:resolve.tex:local-snapshot-20261012; GNU FDL 1.2 or later}
\begin{document}
\maketitle
\noindent\textit{Editorial note.} All delivered mathematical text is the author's own, copied from the pinned Stacks Project chapter file resolve.tex; no statement or proof step was rewritten, reordered or removed.\par\medskip
\noindent\textit{Editorial note.} The counted proof is the author's complete argument: the descent of normality along the faithfully flat completion showing that the source is normal away from the special fibre, the construction through finiteness of normalization over the Nagata completion and the formal glueing lemmas, and the verification of normality at points of the special fibre identifying the constructed finite morphism with the normalization.\par\medskip
\noindent\textit{Editorial note.} The label lemma-normalization-completion names the counted statement, and the proof environment immediately following it in the source is delivered verbatim as its proof.\par\medskip
\noindent\textit{Editorial note.} The counted statement is carried by the author's own lemma environment in the chapter source.\par\medskip
\noindent\textit{Editorial note.} Pre-open gate: python operations/agent-authoring/tools/precard\_gate.py --source resolve.tex --label lemma-normalization-completion --start 3189 --end 3233 --local-id 28088 returned RESULT=OK (L28088 free, occupancy FREE, public interval clear). Pre-assembly tag self-check: canonical tag resolved and not present in reference/public-items.jsonl.\par\medskip
\noindent\textit{Editorial note.} The statement and proof are copied from the pinned chapter file; the cited inputs are kept as external pointers with their canonical Stacks tags.\par\medskip
\noindent\textit{Allowed standard prerequisites (editorial).}\begin{itemize}
\item Descent of normality along a faithfully flat map. Algebra, the lemma that normality descends along a faithfully flat ring map. Supplies the conclusion that the source is normal away from the special fibre, and later the normality of the constructed scheme at points of the special fibre.
\item Finiteness of normalization for Nagata rings and the formal glueing of finite algebras. Algebra, the lemma that a complete Noetherian local ring is Nagata; Morphisms, the lemma that normalization of a Nagata scheme is finite; More on Algebra, the formal glueing lemmas producing a finite algebra over the original ring from one over the completion. Supply the finiteness of the normalization of the base change and the construction of a finite algebra over the original ring recovering it after completion.
\end{itemize}\medskip
\ifdefined\hypertarget\hypertarget{sourcepointer-1}{}\fi
\noindent\textit{Source reference.} Original lemma 032W: sources/stacks-chow/tags.txt, source lines 4035--4035. Prerequisite pointer invoked at proof line(s) 3213; it names a standard fact of a sibling chapter that the counted argument uses, and every step of the counted proof built on it is retained verbatim in the delivered body.\par\medskip
\ifdefined\hypertarget\hypertarget{sourcepointer-2}{}\fi
\noindent\textit{Source reference.} Original lemma 033G: sources/stacks-chow/tags.txt, source lines 4055--4055. Prerequisite pointer invoked at proof line(s) 3208; it names a standard fact of a sibling chapter that the counted argument uses, and every step of the counted proof built on it is retained verbatim in the delivered body.\par\medskip
\ifdefined\hypertarget\hypertarget{sourcepointer-3}{}\fi
\noindent\textit{Source reference.} Original lemma 0ALK: sources/stacks-chow/tags.txt, source lines 13373--13373. Prerequisite pointer invoked at proof line(s) 3222; it names a standard fact of a sibling chapter that the counted argument uses, and every step of the counted proof built on it is retained verbatim in the delivered body.\par\medskip
\ifdefined\hypertarget\hypertarget{sourcepointer-4}{}\fi
\noindent\textit{Source reference.} Original remark 05EU: sources/stacks-chow/tags.txt, source lines 7003--7003. Prerequisite pointer invoked at proof line(s) 3223; it names a standard fact of a sibling chapter that the counted argument uses, and every step of the counted proof built on it is retained verbatim in the delivered body.\par\medskip
\ifdefined\hypertarget\hypertarget{sourcepointer-5}{}\fi
\noindent\textit{Source reference.} Original lemma 035S: sources/stacks-chow/tags.txt, source lines 4139--4139. Prerequisite pointer invoked at proof line(s) 3215; it names a standard fact of a sibling chapter that the counted argument uses, and every step of the counted proof built on it is retained verbatim in the delivered body.\par\medskip
\begin{lemma}
\label{lemma-normalization-completion}
Let $(A, \mathfrak m, \kappa)$ be a Noetherian local ring.
Let $X \to \Spec(A)$ be a morphism which is locally of finite type.
Set $Y = X \times_{\Spec(A)} \Spec(A^\wedge)$.
If the complement of the special fibre in $Y$ is normal, then
the normalization $X^\nu \to X$ is finite and the base change
of $X^\nu$ to $\Spec(A^\wedge)$ recovers the normalization of $Y$.
\end{lemma}

\begin{proof}
There is an immediate reduction to the case where $X = \Spec(B)$
is affine with $B$ a finite type $A$-algebra. Set $C = B \otimes_A A^\wedge$
so that $Y = \Spec(C)$. Since
$A \to A^\wedge$ is faithfully flat, for any prime $\mathfrak q \subset B$
there exists a prime $\mathfrak r \subset C$ lying over $\mathfrak q$.
Then $B_\mathfrak q \to C_\mathfrak r$ is faithfully flat. Hence if
$\mathfrak q$ does not lie over $\mathfrak m$, then $C_\mathfrak r$
is normal by assumption on $Y$ and we conclude that $B_\mathfrak q$
is normal by Algebra, Lemma \ref{algebra-lemma-descent-normal}.
In this way we see that $X$ is normal away from the special fibre.

\medskip\noindent
Recall that the complete Noetherian local ring $A^\wedge$ is Nagata
(Algebra, Lemma \ref{algebra-lemma-Noetherian-complete-local-Nagata}).
Hence the normalization $Y^\nu \to Y$ is finite
(Morphisms, Lemma \ref{morphisms-lemma-nagata-normalization})
and an isomorphism away from the special fibre. Say $Y^\nu = \Spec(C')$.
Then $C \to C'$ is finite and an isomorphism away from $V(\mathfrak m C)$.
Since $B \to C$ is flat and induces an isomorphism
$B/\mathfrak m B \to C/\mathfrak m C$ there exists a finite
ring map $B \to B'$ whose base change to $C$ recovers $C \to C'$.
See More on Algebra, Lemma
\ref{more-algebra-lemma-application-formal-glueing} and
Remark \ref{more-algebra-remark-formal-glueing-algebras}.
Thus we find a finite morphism $X' \to X$ which is an isomorphism
away from the special fibre and whose base change recovers $Y^\nu \to Y$.
By the discussion in the first paragraph we see that $X'$ is normal at
points not on the special fibre. For a point $x \in X'$ on the special
fibre we have a corresponding point $y \in Y^\nu$ and a flat map
$\mathcal{O}_{X', x} \to \mathcal{O}_{Y^\nu, y}$.
Since $\mathcal{O}_{Y^\nu, y}$ is normal, so is $\mathcal{O}_{X', x}$, see
Algebra, Lemma \ref{algebra-lemma-descent-normal}.
Thus $X'$ is normal and it follows that it is the normalization of $X$.
\end{proof}

\end{document}
